% ═══════════════════════════════════════════════════════════════════════
% Higher-Order Residue Deficits in Legendre Intervals
% Titan Project — Paper XII (revised) — February 2026
% ═══════════════════════════════════════════════════════════════════════

\documentclass[11pt, a4paper]{article}

\usepackage[top=28mm, bottom=28mm, left=25mm, right=25mm]{geometry}
\usepackage[T1]{fontenc}
\usepackage{amsmath, amssymb, amsthm, mathtools}
\usepackage[dvipsnames]{xcolor}
\usepackage{enumitem}
\usepackage{booktabs}
\usepackage{hyperref}
\usepackage{float}

\newtheorem{theorem}{Theorem}[section]
\newtheorem{lemma}[theorem]{Lemma}
\newtheorem{corollary}[theorem]{Corollary}
\newtheorem{proposition}[theorem]{Proposition}
\theoremstyle{definition}
\newtheorem{definition}[theorem]{Definition}
\newtheorem{example}[theorem]{Example}
\newtheorem{question}[theorem]{Open Question}
\theoremstyle{remark}
\newtheorem{remark}[theorem]{Remark}

\newcommand{\Leg}[2]{\left(\frac{#1}{#2}\right)}

\title{\textbf{Higher-Order Residue Deficits \\
in Legendre Intervals}}
\author{Ruqing Chen\\[4pt]
\small GUT Geoservice Inc., Montr\'eal, Canada\\[2pt]
\small \texttt{ruqing@hotmail.com}\\[2pt]
\small Repository:
\url{https://github.com/Ruqing1963/higher-order-residue-deficits}}
\date{February 2026}

\begin{document}

\maketitle

\begin{abstract}
We generalize the Spin Asymmetry Theorem of \cite{ChenSpin}
from quadratic residues to arbitrary $k$-th power residues.
Let $n \ge 2$ with $p = 2n - 1$ prime, and let $k \ge 2$
with $k \mid (p-1)$.
The interior of the Legendre interval
$\mathcal{I}_n = [(n-1)^2, n^2]$
exhibits a \emph{universal phase deficit}:
the $k$-th power residue class containing the
missing residue $r \equiv 4^{-1} \pmod{p}$
has exactly one fewer representative than each of the
other $k - 1$ classes.
We identify the deficient phase explicitly:
for any character $\chi$ of order $k$,
$\sum_{x \in \mathcal{I}_n^\circ} \chi(x) = -\chi(4^{-1})$.
For cubic residues ($k = 3$),
the deficient phase is determined by the cubic
character of $2$ modulo $p$.
For quartic residues ($k = 4$),
the deficient phase is restricted to $\Phi_0$ or $\Phi_2$
(never $\Phi_1$ or $\Phi_3$), governed by the
Legendre symbol $\Leg{2}{p}$, hence by $p \bmod 8$.
\end{abstract}

\bigskip

% ═══════════════════════════════════════════════════════════════════════
\section{Introduction}\label{sec:intro}
% ═══════════════════════════════════════════════════════════════════════

In \cite{ChenSpin}, we proved that
for $p = 2n - 1$ prime, the interior
$\mathcal{I}_n^\circ = \{(n-1)^2 + 1, \ldots, n^2 - 1\}$
is a punctured complete residue system modulo $p$,
covering every class except $r \equiv n^2 \pmod{p}$.
Since $r$ is a quadratic residue, $N^- - N^+ = 1$.

At the end of \cite{ChenSpin}, we asked whether
analogous asymmetries hold for higher-order residue symbols.
In this note, we answer affirmatively and go further:
we identify the missing residue as
$r \equiv 4^{-1} \pmod{p}$ and use this to determine
the deficient phase \emph{exactly} for cubic and quartic
residues via classical reciprocity laws.

% ═══════════════════════════════════════════════════════════════════════
\section{The General Theorem}\label{sec:theorem}
% ═══════════════════════════════════════════════════════════════════════

\begin{definition}[$k$-th power residue classes]
\label{def:classes}
Let $p$ be an odd prime and $k \ge 2$ with $k \mid (p-1)$.
Fix a primitive $k$-th root of unity $\zeta \pmod{p}$.
The $k$-th power residue symbol partitions
$(\mathbb{Z}/p\mathbb{Z})^*$ into $k$ classes
$\Phi_j = \{a : a^{(p-1)/k} \equiv \zeta^j \pmod{p}\}$,
each of cardinality $(p-1)/k$.
\end{definition}

\begin{theorem}[Higher-order phase deficit]
\label{thm:main}
Let $n \ge 2$ with $p = 2n - 1$ prime, and let $k \ge 2$
with $k \mid (p-1)$ and $(p-1)/k \ge 2$.
Then the interior $\mathcal{I}_n^\circ$ contains:
\begin{enumerate}[nosep, label=(\roman*)]
\item exactly $(p-1)/k - 1$ elements in the class $\Phi_j$
containing $r \equiv 4^{-1} \pmod{p}$
\textup{(}the deficient phase\textup{)},
\item exactly $(p-1)/k$ elements in each $\Phi_i$
with $i \ne j$,
\item exactly $1$ element divisible by $p$.
\end{enumerate}
\end{theorem}

\begin{proof}
The interval bounds satisfy
$(n-1)^2 = n^2 - (2n - 1) = n^2 - p$,
so $(n-1)^2 \equiv n^2 \pmod{p}$.
Thus $\mathcal{I}_n^\circ$ consists of exactly
$p - 1$ consecutive integers, forming a complete
residue system modulo $p$ punctured at
$r \equiv n^2 \pmod{p}$.

Since $p = 2n - 1$ gives $2n \equiv 1$, we have
$n \equiv 2^{-1}$ and
$r \equiv n^2 \equiv 4^{-1} \pmod{p}$.

Among the $p - 1$ nonzero classes, each phase $\Phi_j$
has $(p-1)/k$ elements.
Removing the single class $r \in \Phi_j$ reduces
$|\Phi_j \cap \mathcal{I}_n^\circ|$ by $1$.
\end{proof}

\begin{corollary}[Character sum identity]
\label{cor:charsum}
For any nontrivial multiplicative character $\chi$
of order dividing $k$ modulo $p$:
\begin{equation}\label{eq:charsum}
    \sum_{x \in \mathcal{I}_n^\circ} \chi(x)
    \;=\; -\chi(r)
    \;=\; -\chi(4^{-1})
    \;=\; -\overline{\chi(2)}^{\,2}.
\end{equation}
This is an exact evaluation.
\end{corollary}

\begin{proof}
$\displaystyle\sum_{x \in \mathcal{I}_n^\circ} \chi(x)
= \sum_{\substack{a=1\\a \ne r}}^{p-1} \chi(a)
= -\chi(r)$,
since $\sum_{a=1}^{p-1} \chi(a) = 0$ for nontrivial
$\chi$.
Then $\chi(4^{-1}) = \overline{\chi(4)}
= \overline{\chi(2)^2} = \overline{\chi(2)}^{\,2}$.
\end{proof}

% ═══════════════════════════════════════════════════════════════════════
\section{Identifying the Deficient Phase}\label{sec:phases}
% ═══════════════════════════════════════════════════════════════════════

\subsection{$k = 2$: Quadratic (Paper IX recovery)}

Since $4^{-1} = (2^{-1})^2$ is a perfect square,
$r \in \Phi_0$ always.
This recovers $N^- - N^+ = 1$.

\subsection{$k = 3$: Cubic residues
($p \equiv 1 \pmod{3}$)}

\begin{proposition}[Cubic deficient phase]
\label{prop:cubic}
The deficient cubic phase is determined by the cubic
character of $2$:
\[
    \chi_3(r) = \chi_3(4^{-1})
    = \overline{\chi_3(2)}^{\,2}
    = \chi_3(2)^4
    = \chi_3(2),
\]
where the last step uses $\chi_3^3 = 1$.
\end{proposition}

Concretely: if $2$ is a cubic residue modulo $p$
($2^{(p-1)/3} \equiv 1$), then the deficit is in $\Phi_0$.
Otherwise, the deficit is in $\Phi_1$ or $\Phi_2$
according to $2^{(p-1)/3} \pmod{p}$.

By classical algebraic number theory, $2$ is a cubic
residue modulo a prime $p \equiv 1 \pmod{3}$ if and
only if $p$ is representable as $p = x^2 + 27y^2$ for
some integers $x, y$.

\subsection{$k = 4$: Quartic residues
($p \equiv 1 \pmod{4}$)}

\begin{proposition}[Quartic phase restriction]
\label{prop:quartic}
The deficient quartic phase is always $\Phi_0$ or $\Phi_2$
\textup{(}never $\Phi_1$ or $\Phi_3$\textup{)}.
Specifically:
\[
    \chi_4(r)
    = \chi_4(4^{-1})
    = \chi_4(2)^{-2}
    = \chi_2(2)^{-1}
    = \Leg{2}{p},
\]
since $\chi_4^2 = \chi_2$.
Therefore:
\begin{itemize}[nosep]
\item $p \equiv 1 \pmod{8}$: $\Leg{2}{p} = 1$,
deficit in $\Phi_0$.
\item $p \equiv 5 \pmod{8}$: $\Leg{2}{p} = -1$,
deficit in $\Phi_2$.
\end{itemize}
\end{proposition}

\begin{proof}
$\chi_4(4^{-1})
= \chi_4(2^{-2})
= \chi_4(2)^{-2}$.
Since $\chi_4^2 = \chi_2$,
$\chi_4(2)^{-2} = \chi_2(2)^{-1} = \Leg{2}{p}$,
which equals $\pm 1$.
Since $\Phi_0$ corresponds to character value $1$ and
$\Phi_2$ to $-1$, the deficit is confined to these two.
\end{proof}

% ═══════════════════════════════════════════════════════════════════════
\section{Computational Illustration}\label{sec:computation}
% ═══════════════════════════════════════════════════════════════════════

Since Theorem~\ref{thm:main} is an unconditional algebraic
identity, its computational verification is a tautology.
We display sample data to illustrate the mechanism.

\begin{table}[H]
\centering
\caption{Phase deficit verification for $k = 2, 3, 4, 5$,
all qualifying primes up to $n = 1000$.}
\label{tab:verification}
\medskip
\begin{tabular}{@{}c r c@{}}
\toprule
$k$ & Qualifying primes & Deficit confirmed \\
\midrule
2 (quadratic) & 549 & $\checkmark$ (all) \\
3 (cubic)     & 148 & $\checkmark$ (all) \\
4 (quartic)   & 146 & $\checkmark$ (all) \\
5 (quintic)   & 73  & $\checkmark$ (all) \\
\bottomrule
\end{tabular}
\end{table}

\begin{table}[H]
\centering
\caption{Cubic phase deficit ($k = 3$).
$*$ marks the deficient class.
The final column shows $2^{(p-1)/3} \bmod p$.}
\label{tab:cubic}
\medskip
\begin{tabular}{@{}r r r r r c@{}}
\toprule
$n$ & $p$ & $|\Phi_0|$ & $|\Phi_1|$ & $|\Phi_2|$
& $2^{(p-1)/3}$ \\
\midrule
 7 & 13 &  4  & $3^*$ &  4 & $\omega$ \\
10 & 19 &  6  & $5^*$ &  6 & $\omega$ \\
16 & 31 & $9^*$ & 10  & 10 & $1$ \\
22 & 43 & $13^*$ & 14 & 14 & $1$ \\
34 & 67 & 22 & 22 & $21^*$ & $\omega^2$ \\
\bottomrule
\end{tabular}
\end{table}

\begin{table}[H]
\centering
\caption{Quartic phase restriction ($k = 4$).
The deficit is always in $\Phi_0$ or $\Phi_2$,
governed by $p \bmod 8$.}
\label{tab:quartic}
\medskip
\begin{tabular}{@{}r r c r r r r c@{}}
\toprule
$n$ & $p$ & $p\!\bmod 8$ &
$|\Phi_0|$ & $|\Phi_1|$ & $|\Phi_2|$ & $|\Phi_3|$ &
Deficit \\
\midrule
 9 & 17 & 1 & $3^*$ & 4 & 4 & 4 & $\Phi_0$ \\
15 & 29 & 5 & 7 & 7 & $6^*$ & 7 & $\Phi_2$ \\
21 & 41 & 1 & $9^*$ & 10 & 10 & 10 & $\Phi_0$ \\
27 & 53 & 5 & 13 & 13 & $12^*$ & 13 & $\Phi_2$ \\
\bottomrule
\end{tabular}
\end{table}

% ═══════════════════════════════════════════════════════════════════════
\section{Discussion}\label{sec:discussion}
% ═══════════════════════════════════════════════════════════════════════

\subsection{The role of $4^{-1}$}

The identification $r \equiv 4^{-1} \pmod{p}$ is the
key algebraic insight.
It reduces the study of the deficient phase to the
classical problem of determining the $k$-th power
residue character of $2$ modulo $p$.
For $k = 2$: trivially $\Phi_0$ (since $4^{-1}$ is a square).
For $k = 3$: connects to cubic reciprocity and
the representation $p = x^2 + 27y^2$.
For $k = 4$: reduces to the Legendre symbol
$\Leg{2}{p}$ via the supplementary law.

\subsection{Simultaneous deficits}

When $p - 1$ is highly composite, the interval
simultaneously exhibits compatible deficits across
all dividing $k$, all from the single puncture
at $4^{-1}$.

\subsection{Open questions}

\begin{question}[Cubic Oppermann analogue]
For $k = 3$, does the right half of $\mathcal{I}_n$
exhibit zero cubic skewness under some congruence
condition on $n$?
\end{question}

\begin{question}[Beyond Burgess for Legendre intervals]
The identity \eqref{eq:charsum} evaluates the full
character sum exactly.
For half-intervals, Burgess \cite{Burgess1963} gives
$O(p^{3/16 + \varepsilon})$ bounds for short sums.
Can the arithmetic structure of Legendre half-intervals
yield improvements for these specific ranges?
\end{question}

\subsection*{Data and code availability}

Verification scripts are available at
\url{https://github.com/Ruqing1963/higher-order-residue-deficits}.

% ═══════════════════════════════════════════════════════════════════════
\begin{thebibliography}{99}

\bibitem{ChenSpin}
  R.~Chen,
  ``Algebraic rigidity and quadratic residue asymmetry
  in Legendre intervals,''
  Zenodo, 2026.
  \url{https://zenodo.org/records/18706876}

\bibitem{ChenOpp}
  R.~Chen,
  ``Oppermann's parity law: quadratic residue symmetry
  breaking in half-intervals via the negation involution,''
  Zenodo, 2026.
  \url{https://zenodo.org/records/18707265}

\bibitem{Burgess1963}
  D.\,A.~Burgess,
  ``On character sums and $L$-series.\ II,''
  \textit{Proc.\ London Math.\ Soc.}\ (3) \textbf{13}
  (1963), 524--536.

\end{thebibliography}

\end{document}
